% Computer Science A --- Static front-end (HTML, CSS).

\runningheader{Computer Science A}{Web: HTML \& CSS}{Page \thepage\ of \numpages}

\begingradingrange{csawebstatic}

\uplevel{\par\textbf{Part 1: HTML \& CSS basics}\par}

\question[4]
What does HTML (HyperText Markup Language) primarily define for a web page?
\begin{choices}
  \choice The visual theme and colors only
  \choice Only interactive behavior and event handlers
  \CorrectChoice The structure and semantic content of the document
  \choice The execution order of JavaScript modules
\end{choices}
\begin{solution}
  \textbf{HTML} marks up headings, paragraphs, lists, links, and other content so the browser can build the DOM tree. \textbf{CSS} handles presentation; \textbf{JavaScript} adds behavior.
\begin{lstlisting}[style=htmlexam]
<!DOCTYPE html>
<html>
  <body>
    <h1>Title</h1>
    <p>Paragraphs and lists carry meaning, not just looks.</p>
  </body>
</html>
\end{lstlisting}
\end{solution}

\question[4]
Which element sets the text typically shown in the browser tab / window title?
\begin{choices}
  \choice \texttt{<h1>}
  \choice \texttt{<header>}
  \CorrectChoice \texttt{<title>}
  \choice \texttt{<meta>}
\end{choices}
\begin{solution}
  The \texttt{<title>} element lives in \texttt{<head>} and names the document. Headings \texttt{<h1>}--\texttt{<h6>} structure visible page content.
\begin{lstlisting}[style=htmlexam]
<head>
  <meta charset="utf-8">
  <title>My lab page</title>
</head>
\end{lstlisting}
\end{solution}

\question[4]
Which fragment correctly represents an \textbf{unordered} list with two items?
\begin{choices}
  \choice \texttt{<ol><li>A</li><li>B</li></ol>}
  \CorrectChoice \texttt{<ul><li>A</li><li>B</li></ul>}
  \choice \texttt{<list><item>A</item></list>}
  \choice \texttt{<ul>A</ul><ul>B</ul>}
\end{choices}
\begin{solution}
  Use \texttt{<ul>} with nested \texttt{<li>} items for bullet lists; \texttt{<ol>} is for numbered lists. Each list item must be an \texttt{<li>}.
\begin{lstlisting}[style=htmlexam]
<ul>
  <li>A</li>
  <li>B</li>
</ul>
\end{lstlisting}
\end{solution}

\question[4]
In HTML, which pair correctly matches document metadata and visible page content?
\begin{choices}
  \CorrectChoice \texttt{<head>} for metadata and \texttt{<body>} for visible content
  \choice \texttt{<body>} for metadata and \texttt{<head>} for visible content
  \choice \texttt{<title>} for metadata and \texttt{<main>} for all visible content
  \choice \texttt{<meta>} for metadata and \texttt{<html>} for visible content only
\end{choices}
\begin{solution}
  \texttt{<head>} stores metadata (title, meta tags, stylesheet links), while \texttt{<body>} contains what users see in the rendered page.
\end{solution}

\question[4]
What syntax is HTML primarily made of?
\begin{choices}
  \CorrectChoice Tags (elements and attributes)
  \choice SQL clauses
  \choice Java bytecode
  \choice CSS declarations only
\end{choices}
\begin{solution}
  HTML is markup composed of tags such as \texttt{<p>}, \texttt{<a>}, and \texttt{<div>}, with attributes for additional semantics and behavior.
\end{solution}

\question[4]
What should appear first in a standard HTML5 document?
\begin{choices}
  \choice \texttt{<html>}
  \CorrectChoice \texttt{<!DOCTYPE html>}
  \choice \texttt{<head>}
  \choice \texttt{<meta charset="utf-8">}
\end{choices}
\begin{solution}
  The doctype declaration tells the browser to render in HTML5 standards mode.
\end{solution}

\question[4]
In which parent element(s) are \texttt{<li>} list items valid?
\begin{choices}
  \choice Only \texttt{<ul>}
  \choice Only \texttt{<ol>}
  \CorrectChoice Both \texttt{<ul>} and \texttt{<ol>}
  \choice \texttt{<table>}
\end{choices}
\begin{solution}
  List items belong inside ordered lists (\texttt{<ol>}) or unordered lists (\texttt{<ul>}).
\end{solution}

\question[4]
Which of the following is \textbf{not} a semantic HTML element?
\begin{choices}
  \choice \texttt{<article>}
  \choice \texttt{<nav>}
  \CorrectChoice \texttt{<div>}
  \choice \texttt{<section>}
\end{choices}
\begin{solution}
  \texttt{<div>} is non-semantic and primarily used as a generic container; semantic elements convey structure and meaning.
\end{solution}

\question[4]
Which HTML element is a void element (does not require a closing tag)?
\begin{choices}
  \CorrectChoice \texttt{<img>}
  \choice \texttt{<p>}
  \choice \texttt{<li>}
  \choice \texttt{<table>}
\end{choices}
\begin{solution}
  \texttt{<img>} is a void element. In HTML5, void elements do not wrap content and do not use a separate closing tag.
\end{solution}

\question[4]
Which markup is the most correct table structure?
\begin{choices}
  \CorrectChoice
\begin{lstlisting}[style=htmlexam]
<table>
  <tr><th>Header 1</th><th>Header 2</th></tr>
  <tr><td>Row 1, Cell 1</td><td>Row 1, Cell 2</td></tr>
  <tr><td>Row 2, Cell 1</td><td>Row 2, Cell 2</td></tr>
</table>
\end{lstlisting}
  \choice
\begin{lstlisting}[style=htmlexam]
<table>
  <tr><th>Header 1</th><th>Header 2
  <tr><td>Row 1, Cell 1</td><td>Row 1, Cell 2</td></tr>
  <td>Row 2, Cell 1</td><td>Row 2, Cell 2</td>
</table>
\end{lstlisting}
  \choice
\begin{lstlisting}[style=htmlexam]
<table>
  <thead>Header 1 Header 2</thead>
  <tbody>Row 1, Cell 1 Row 1, Cell 2</tbody>
</table>
\end{lstlisting}
  \choice
\begin{lstlisting}[style=htmlexam]
<table>
  <td>Header 1</td><td>Header 2</td>
</table>
\end{lstlisting}
\end{choices}
\begin{solution}
  Proper tables use row containers (\texttt{<tr>}) and then header/data cells (\texttt{<th>} / \texttt{<td>}) within each row.
\end{solution}

\question[4]
Which syntax is correct for an image element with source and alternate text?
\begin{choices}
  \CorrectChoice \texttt{<img src="image.jpg" alt="Description">}
  \choice \texttt{<img href="image.jpg" alt="Description"></img>}
  \choice \texttt{<image src="image.jpg" alt="Description">}
  \choice \texttt{<img>image.jpg</img>}
\end{choices}
\begin{solution}
  Use \texttt{src} for the image URL/path and \texttt{alt} for accessible fallback text.
\end{solution}

\question[4]
Which fragment is a correct minimal HTML5 boilerplate skeleton?
\begin{choices}
  \CorrectChoice
\begin{lstlisting}[style=htmlexam]
<!DOCTYPE html>
<html>
  <head><title>Page Title</title></head>
  <body><!-- Content goes here --></body>
</html>
\end{lstlisting}
  \choice
\begin{lstlisting}[style=htmlexam]
<html>
  <title>Page Title</title>
  <body>Content</body>
</html>
\end{lstlisting}
  \choice
\begin{lstlisting}[style=htmlexam]
<!DOCTYPE html>
<head><body>Content</body></head>
\end{lstlisting}
  \choice
\begin{lstlisting}[style=htmlexam]
<body>
  <head><title>Page Title</title></head>
</body>
\end{lstlisting}
\end{choices}
\begin{solution}
  A correct skeleton starts with HTML5 doctype, then \texttt{<html>}, then separate \texttt{<head>} and \texttt{<body>} sections.
\end{solution}


\newpage

\question[4]
Which element wraps items that should appear as a \textbf{numbered} (ordered) list?
\begin{choices}
  \choice \texttt{<li>}
  \choice \texttt{<dl>}
  \CorrectChoice \texttt{<ol>}
  \choice \texttt{<nav>}
\end{choices}
\begin{solution}
  \texttt{<ol>} (ordered list) works with \texttt{<li>} children for step-like or ranked content; \texttt{<ul>} is for bullets. Definition lists use \texttt{<dl>}/\texttt{<dt>}/\texttt{<dd>}.
\begin{lstlisting}[style=htmlexam]
<ol>
  <li>Step one</li>
  <li>Step two</li>
</ol>
\end{lstlisting}
\end{solution}

\question[4]
In \texttt{<a href="page.html">Open</a>}, what is the role of \texttt{href}?
\begin{choices}
  \CorrectChoice It specifies the link destination (URL or path)
  \choice It sets the font size of the link text
  \choice It makes the link open only in a new window
  \choice It prevents the default navigation behavior
\end{choices}
\begin{solution}
  The \textbf{hypertext reference} attribute \texttt{href} points to the resource to load. Presentation is CSS; \texttt{target} and JS handlers affect navigation behavior separately.
\begin{lstlisting}[style=htmlexam]
<a href="page.html">Open</a>
<a href="https://example.com/" target="_blank">External</a>
\end{lstlisting}
\end{solution}

\question[4]
In the lab progression, \textbf{inline CSS} is applied by placing declarations where?
\begin{choices}
  \choice Only in an external \texttt{.css} file linked via \texttt{<link>}
  \choice Only inside a \texttt{<style>} block in the head
  \CorrectChoice In the element's \texttt{style} attribute as a declaration list
  \choice In a JavaScript \texttt{className} assignment only
\end{choices}
\begin{solution}
  Inline styles use \texttt{style="property: value; ..."} on the element. External and embedded stylesheets are the other main approaches; they cascade with specificity rules.
\begin{lstlisting}[style=htmlexam]
<p style="color: navy; font-size: 14px;">Styled in place</p>
\end{lstlisting}
\end{solution}

\question[4]
Which option shows the most correct syntax for a basic login-style HTML form?
\begin{choices}
  \CorrectChoice
\begin{lstlisting}[style=htmlexam]
<form>
  <label for="username">Username:</label>
  <input id="username" type="text" name="username">
  <label for="password">Password:</label>
  <input id="password" type="password" name="password">
  <input type="submit" value="Submit">
</form>
\end{lstlisting}
  \choice
\begin{lstlisting}[style=htmlexam]
<form>
  <label username>Username</label>
  <input text="username">
</form>
\end{lstlisting}
  \choice
\begin{lstlisting}[style=htmlexam]
<form>
  <input id="username">
  <input id="password">
  <button>
</form>
\end{lstlisting}
  \choice
\begin{lstlisting}[style=htmlexam]
<form action>
  <label>Username</label>
  <input>
  <submit>Submit</submit>
</form>
\end{lstlisting}
\end{choices}
\begin{solution}
  A well-formed example uses explicit \texttt{label}--\texttt{for}/\texttt{id} linkage and valid input \texttt{type} and \texttt{name} attributes.
\end{solution}

\question[4]
Which of the following is \textbf{not} part of the CSS box model?
\begin{choices}
  \choice Content
  \choice Padding
  \choice Border
  \CorrectChoice Whitespace
\end{choices}
\begin{solution}
  The canonical box model layers are content, padding, border, and margin.
\end{solution}

\question[4]
Which style source has the highest normal cascade precedence (ignoring \texttt{!important})?
\begin{choices}
  \CorrectChoice Inline styles (style attribute)
  \choice Embedded stylesheet rules only
  \choice External stylesheet rules only
  \choice Browser default stylesheet
\end{choices}
\begin{solution}
  Inline styles are typically strongest in normal cascade order, then selector specificity and source order resolve ties among stylesheet rules.
\end{solution}

\question[4]
In CSS, which selector targets a class?
\begin{choices}
  \CorrectChoice \texttt{.myClass}
  \choice \texttt{\#myClass}
  \choice \texttt{myClass}
  \choice \texttt{*myClass}
\end{choices}
\begin{solution}
  Class selectors begin with a period; ID selectors begin with \texttt{\#}.
\end{solution}

\question[4]
In CSS, which selector targets an element with a specific id?
\begin{choices}
  \choice \texttt{.myId}
  \CorrectChoice \texttt{\#myId}
  \choice \texttt{myId()}
  \choice \texttt{*myId}
\end{choices}
\begin{solution}
  The hash prefix identifies ID selectors.
\end{solution}

\question[4]
Given \texttt{.container \{ background-color: lightblue; \}}, what is the div background color?
\begin{choices}
  \CorrectChoice Light blue
  \choice Green
  \choice Transparent only
  \choice Browser default gray
\end{choices}
\begin{solution}
  The declaration explicitly sets the rendered background to light blue.
\end{solution}

\question[4]
Given \texttt{.text \{ color: green; font-size: 16px; \}}, what is the paragraph text color?
\begin{choices}
  \choice Blue
  \CorrectChoice Green
  \choice Black
  \choice Inherited and therefore unknown
\end{choices}
\begin{solution}
  \texttt{color: green;} directly sets the text color to green.
\end{solution}

\question[4]
Given \texttt{.center-text \{ text-align: center; \}}, how is the paragraph text aligned?
\begin{choices}
  \choice Left
  \CorrectChoice Center
  \choice Right
  \choice Justified
\end{choices}
\begin{solution}
  \texttt{text-align: center;} centers inline text content within the element's box.
\end{solution}

\question[4]
Given \texttt{.box \{ border: 2px solid red; \}}, what is the border color?
\begin{choices}
  \choice Blue
  \choice Black
  \CorrectChoice Red
  \choice Transparent
\end{choices}
\begin{solution}
  The shorthand \texttt{border} sets width, style, and color in one declaration; here the color token is red.
\end{solution}

\question[4]
Given the snippet with a yellow container and inner text styled as blue with margin/padding/red border, which description is most accurate?
\begin{choices}
  \CorrectChoice Yellow container background with blue text; the text element also has 20px margin, 10px padding, and a 2px red border
  \choice Blue container background with yellow text and no spacing
  \choice Yellow text with blue border on the container
  \choice No visible style differences because classes do not apply to nested elements
\end{choices}
\begin{solution}
  The outer class styles the container's dimensions and yellow background, while the inner class styles text color and spacing/border around that paragraph element.
\end{solution}


\newpage

\question[4]
Which CSS declaration horizontally centers text inside its container (for typical block contexts in the lab)?
\begin{choices}
  \choice \texttt{font-style: italic;}
  \choice \texttt{font-size: 14px;}
  \CorrectChoice \texttt{text-align: center;}
  \choice \texttt{margin: 0 auto;}
\end{choices}
\begin{solution}
  \texttt{text-align} affects inline content within a block. Centering a block itself often uses margins or flexbox; the Homepage/CSS labs emphasized \texttt{text-align} for typography alignment.
\begin{lstlisting}[style=cssexam]
h1 {
  text-align: center;
}
\end{lstlisting}
\end{solution}

\question[4]
The CSS Inline lab emphasizes \texttt{font-style}, \texttt{font-size}, and \texttt{text-align}. Which value of \texttt{font-style} produces italicized text?
\begin{choices}
  \choice \texttt{bold}
  \CorrectChoice \texttt{italic}
  \choice \texttt{underline}
  \choice \texttt{sans-serif}
\end{choices}
\begin{solution}
  \texttt{font-style: italic} slants glyphs; \texttt{font-size} sets length (e.g.\ \texttt{14px}). \texttt{bold} belongs to \texttt{font-weight}; \texttt{underline} is usually \texttt{text-decoration}.
\begin{lstlisting}[style=cssexam]
.note {
  font-style: italic;
  font-size: 14px;
}
\end{lstlisting}
\end{solution}

\question[4]
The rule \texttt{h1 \{ color: blue; \}} applies to which elements?
\begin{choices}
  \CorrectChoice Every \texttt{h1} element (type / element selector)
  \choice Only \texttt{h1} elements with \texttt{id="color"}
  \choice Only \texttt{h1} elements whose \texttt{class} is \texttt{blue}
  \choice The first line of every paragraph
\end{choices}
\begin{solution}
  A bare tag name is an \textbf{element} (type) selector. Combined with classes or IDs it gains specificity; the CSS Selectors lab contrasted element, \texttt{.class}, and \texttt{\#id}.
\begin{lstlisting}[style=cssexam]
h1 {
  color: blue;
}
\end{lstlisting}
\end{solution}

\newpage


\question[4]
In CSS, the selector \texttt{.highlight} targets which elements?
\begin{choices}
  \choice Every element with tag name \texttt{highlight}
  \CorrectChoice Every element whose \texttt{class} attribute includes \texttt{highlight}
  \choice The unique element with \texttt{id="highlight"}
  \choice Only \texttt{<span>} elements
\end{choices}
\begin{solution}
  The \textbf{class selector} prefix \texttt{.} matches the space-separated \texttt{class} list. \textbf{ID} uses \texttt{\#}; \textbf{type} selectors use the bare tag name (\texttt{h1}, etc.).
\begin{lstlisting}[style=cssexam]
.highlight {
  background-color: yellow;
}
\end{lstlisting}
\begin{lstlisting}[style=htmlexam]
<p class="highlight">Important</p>
\end{lstlisting}
\end{solution}

\question[4]
In the CSS box model, which area is \textbf{outside} the border (separating this element from neighbors)?
\begin{choices}
  \choice Padding
  \CorrectChoice Margin
  \choice Content box only
  \choice \texttt{outline} (always identical to border)
\end{choices}
\begin{solution}
  From inside out: \textbf{content} $\rightarrow$ \textbf{padding} $\rightarrow$ \textbf{border} $\rightarrow$ \textbf{margin}. Padding is inside the border; margin is outside.
\begin{lstlisting}[style=cssexam]
.box {
  padding: 8px;
  border: 1px solid #333;
  margin: 16px; /* outside the border */
}
\end{lstlisting}
\end{solution}

\question[4]
In the Box Model lab, \texttt{background-color} applied to an element typically fills which regions (in the standard \texttt{box-sizing: content-box} mental model)?
\begin{choices}
  \choice Only the margin
  \choice Only the content box, excluding padding
  \CorrectChoice Content and padding (under the border edge)
  \choice Only areas outside the border
\end{choices}
\begin{solution}
  Background painting defaults to \textbf{content + padding} up to the inner edge of the border (details vary slightly by shorthand and clipping). The lab tied \texttt{background-color} to the boxed area students style alongside \texttt{padding}, \texttt{margin}, and \texttt{border}.
\begin{lstlisting}[style=cssexam]
.panel {
  padding: 12px;
  background-color: #e8f4ff;
}
\end{lstlisting}
\end{solution}


\newpage
\question[4]
When two rules conflict and both apply to the same element, a rule whose selector is only an \textbf{ID} (\texttt{\#main-title}) will usually beat a rule whose selector is only a \textbf{class} (\texttt{.subtitle}) because:
\begin{choices}
  \choice IDs load later in the file
  \CorrectChoice ID selectors have higher specificity weight than a single class selector
  \choice Classes are deprecated in modern CSS
  \choice The browser always prefers the shorter rule
\end{choices}
\begin{solution}
  Specificity approximates \textbf{inline} $>$ \textbf{ID} $>$ \textbf{class/attribute} $>$ \textbf{type}. Tie-breaking uses source order and \texttt{!important} (sparingly). The lab table summarized ID vs class vs element priority.
\begin{lstlisting}[style=cssexam]
.subtitle { color: gray; }
#main-title { color: navy; /* wins over .subtitle */ }
\end{lstlisting}
\end{solution}

\question[4]
The line \texttt{<meta charset="utf-8">} in \texttt{<head>} primarily tells the browser:
\begin{choices}
  \choice Which JavaScript file to load first
  \choice The exact width of the layout viewport
  \CorrectChoice How to interpret the document's bytes as characters (text encoding)
  \choice That the page must not be cached
\end{choices}
\begin{solution}
  The charset declaration names an encoding (here UTF-8) so the parser maps downloaded bytes to the right Unicode characters---especially important for accents and symbols.
\begin{lstlisting}[style=htmlexam]
<head>
  <meta charset="utf-8">
  <title>Encoding demo</title>
</head>
\end{lstlisting}
\end{solution}

\question[4]
Which markup correctly links an \textbf{external} stylesheet \texttt{styles.css} from the same folder?
\begin{choices}
  \choice \texttt{<script href="styles.css"></script>}
  \CorrectChoice \texttt{<link rel="stylesheet" href="styles.css">}
  \choice \texttt{<style src="styles.css"></style>}
  \choice \texttt{<css href="styles.css" />}
\end{choices}
\begin{solution}
  \texttt{<link rel="stylesheet" href="...">} pulls a \texttt{.css} file into the cascade. \texttt{<style>} holds inline CSS text; \texttt{<script>} is for scripts.
\begin{lstlisting}[style=htmlexam]
<head>
  <link rel="stylesheet" href="styles.css">
</head>
\end{lstlisting}
\end{solution}

\newpage

\question[4]
In HTML, \texttt{class="card featured"} on one element means:
\begin{choices}
  \choice The element has invalid markup (only one class name is allowed)
  \choice \texttt{featured} is a tag name nested inside \texttt{card}
  \CorrectChoice The element has \textbf{both} classes \texttt{card} and \texttt{featured} (space-separated list)
  \choice \texttt{card} is the ID and \texttt{featured} is the tag
\end{choices}
\begin{solution}
  The \texttt{class} attribute is a space-separated token list. Selectors \texttt{.card}, \texttt{.featured}, and \texttt{.card.featured} can all match the same node.
\begin{lstlisting}[style=htmlexam]
<article class="card featured">Sale item</article>
\end{lstlisting}
\begin{lstlisting}[style=cssexam]
.card { border: 1px solid #ccc; }
.featured { background: #fff8e6; }
\end{lstlisting}
\end{solution}

\question[4]
In CSS, the selector \texttt{*} (asterisk alone) means:
\begin{choices}
  \choice Only elements named \texttt{star}
  \choice Only the first child of every parent
  \CorrectChoice All elements in the document (universal selector)
  \choice A syntax error in modern CSS
\end{choices}
\begin{solution}
  The universal selector matches every element; it is often used in resets or broad rules (use carefully because it can hurt performance and specificity intuition).
\begin{lstlisting}[style=cssexam]
* { box-sizing: border-box; }
\end{lstlisting}
\end{solution}

\question[4]
Without extra rules, which property is typically \textbf{inherited} from a parent element to its descendants?
\begin{choices}
  \choice \texttt{margin-top}
  \choice \texttt{border-width}
  \CorrectChoice \texttt{color}
  \choice \texttt{width}
\end{choices}
\begin{solution}
  Many text-oriented properties (\texttt{color}, \texttt{font-family}, \texttt{font-size} in the normal case) inherit so children match the parent's typography. Box properties like \texttt{margin} and \texttt{border} usually do not inherit.
\begin{lstlisting}[style=htmlexam]
<section style="color: navy;">
  <p>Inherited text color.</p>
</section>
\end{lstlisting}
\end{solution}

\newpage

\question[4]
The \texttt{alt} attribute on \texttt{<img>} is important mainly because it:
\begin{choices}
  \choice Sets the image file path (same role as \texttt{src})
  \choice Forces the image to load faster
  \CorrectChoice Provides a text substitute for the image (accessibility, broken image, screen readers)
  \choice Defines the clickable region for the image only
\end{choices}
\begin{solution}
  \texttt{alt} describes the image content for assistive tech and when the asset fails to load; \texttt{src} is the actual URL/path.
\begin{lstlisting}[style=htmlexam]
<img src="logo.png" alt="Company logo">
\end{lstlisting}
\end{solution}

\question[4]
In HTML5, \texttt{<br>} and \texttt{<img>} are usually written without a closing tag because they are:
\begin{choices}
  \choice Deprecated and should be replaced by \texttt{<div>}
  \choice Only valid inside \texttt{<table>}
  \CorrectChoice Void (empty) elements---they do not wrap text content
  \choice Required to use self-closing XML form \texttt{</img>} in all browsers
\end{choices}
\begin{solution}
  Void elements have no element content; the slash-only XHTML-style \texttt{<br />} is tolerated but optional in HTML5 for these tags.
\begin{lstlisting}[style=htmlexam]
<p>Line one<br>Line two</p>
<img src="pic.png" alt="">
\end{lstlisting}
\end{solution}

\question[4]
Which statement best describes \textbf{semantic} HTML elements?
\begin{choices}
  \choice A generic \texttt{<div>} is the main example of a semantic element
  \choice Semantic tags are unrelated to document structure
  \CorrectChoice They name meaningful regions (e.g.\ \texttt{<header>}, \texttt{<nav>}, \texttt{<article>}) so structure and grouping carry meaning
  \choice They are developer-only comments that browsers ignore
\end{choices}
\begin{solution}
  Semantic elements describe roles (navigation, main content, aside) rather than being purely presentational wrappers. \texttt{<div>} is non-semantic until given a role or ARIA.
\end{solution}

\question[4]
What is the HTML5 doctype declaration?
\begin{choices}
  \CorrectChoice \texttt{<!DOCTYPE html>}
  \choice \texttt{<!DOCTYPE HTML PUBLIC "-//W3C//DTD HTML 4.01//EN" "http://www.w3.org/TR/html4/strict.dtd">}
  \choice \texttt{<!HTML doctype>}
  \choice \texttt{<DOCTYPE html>}
\end{choices}
\begin{solution}
  HTML5 uses the short, case-insensitive \texttt{<!DOCTYPE html>} so browsers use standards mode.
\begin{lstlisting}[style=htmlexam]
<!DOCTYPE html>
<html lang="en">
  <head><meta charset="utf-8"><title>Demo</title></head>
  <body></body>
</html>
\end{lstlisting}
\end{solution}

\question[4]
Which of these is \textbf{not} a valid CSS unit?
\begin{choices}
  \choice \texttt{\%} (percentage)
  \choice \texttt{em} (font-relative)
  \CorrectChoice \texttt{li} (not a unit---it names a list item in HTML)
  \choice \texttt{px} or \texttt{pt}
\end{choices}
\begin{solution}
  Common units include \texttt{px}, \texttt{em}, \texttt{rem}, \texttt{\%}, \texttt{vh}/\texttt{vw}, \texttt{pt}, etc. \texttt{li} is an element, not a length unit.
\end{solution}

\question[4]
HTML is a programming language.
\begin{choices}
  \choice True
  \CorrectChoice False---HTML is a \textbf{markup} language for structure and semantics, not a general-purpose programming language
  \choice True only when scripts are disabled
  \choice True only for static pages
\end{choices}
\begin{solution}
  Behavior and computation belong to languages like JavaScript; HTML declares content and structure.
\end{solution}

\question[4]
Which group lists elements that normally belong inside \texttt{<body>}?
\begin{choices}
  \CorrectChoice \texttt{<div>}, \texttt{<header>}, \texttt{<span>}, \texttt{<strong>}
  \choice \texttt{<link>}, \texttt{<meta>}, \texttt{<p>}, \texttt{<title>}
  \choice \texttt{<head>}, \texttt{<h1>}, \texttt{<h2>}, \texttt{<h3>} (only in \texttt{<head>})
  \choice \texttt{<article>}, \texttt{<aside>}, \texttt{<nav>}, \texttt{<title>}
\end{choices}
\begin{solution}
  \texttt{<title>}, \texttt{<meta>}, and \texttt{<link>} (for metadata/stylesheets) live in \texttt{<head>}. Visible flow content usually sits in \texttt{<body>}.
\end{solution}

\question[4]
Which values are \textbf{valid} for an \texttt{<input>} element's \texttt{type} attribute?
\begin{choices}
  \CorrectChoice \texttt{password} and \texttt{text} are both valid \texttt{type} values
  \choice \texttt{title} and \texttt{name} are valid \texttt{type} values
  \choice Only \texttt{name} is valid; \texttt{text} is not
  \choice Only \texttt{password} exists in HTML5
\end{choices}
\begin{solution}
  \texttt{name} and \texttt{title} are attributes, not \texttt{type} keywords. Many input types exist (\texttt{text}, \texttt{password}, \texttt{email}, \texttt{checkbox}, etc.).
\begin{lstlisting}[style=htmlexam]
<input type="text" name="user">
<input type="password" name="pass">
\end{lstlisting}
\end{solution}

\question[4]
Which of the following is \textbf{not} an HTML5 \emph{element} name?
\begin{choices}
  \choice \texttt{<body>}
  \choice \texttt{<h1>}
  \choice \texttt{<p>}
  \CorrectChoice \texttt{<class>} (\texttt{class} is an \emph{attribute}, not a tag)
\end{choices}
\begin{solution}
  Use \texttt{class="..."} on elements; there is no \texttt{<class>} tag.
\end{solution}

\question[4]
Order the CSS box model layers from \textbf{outermost} to \textbf{innermost}.
\begin{choices}
  \choice border $\rightarrow$ margin $\rightarrow$ content $\rightarrow$ padding
  \choice margin $\rightarrow$ padding $\rightarrow$ border $\rightarrow$ content
  \CorrectChoice margin $\rightarrow$ border $\rightarrow$ padding $\rightarrow$ content
  \choice border $\rightarrow$ content $\rightarrow$ margin $\rightarrow$ padding
\end{choices}
\begin{solution}
  From outside in: \textbf{margin}, \textbf{border}, \textbf{padding}, \textbf{content}. This matches the mental model used alongside \texttt{background-color} painting under the border.
\end{solution}

\question[4]
Which element holds an \textbf{embedded} (internal) stylesheet in the document?
\begin{choices}
  \choice \texttt{<script>}
  \choice \texttt{<body>}
  \CorrectChoice \texttt{<style>} (usually in \texttt{<head>})
  \choice \texttt{<css>}
\end{choices}
\begin{solution}
  Put CSS rules inside \texttt{<style>}; link external files with \texttt{<link rel="stylesheet" href="...">}.
\begin{lstlisting}[style=htmlexam]
<head>
  <style>
    body { font-family: system-ui, sans-serif; }
  </style>
</head>
\end{lstlisting}
\end{solution}

\question[4]
What is true of a typical \textbf{inline} element (e.g.\ \texttt{<span>}) in normal flow?
\begin{choices}
  \CorrectChoice It sits in the text line by default; horizontal size follows content (subject to wrapping)
  \choice It always starts on a new line like a block box
  \choice It always stretches to the full viewport width
  \choice The browser truncates its text when width is tight
\end{choices}
\begin{solution}
  Inline boxes participate in line layout; block-level boxes (e.g.\ \texttt{<div>}) typically break onto new lines and can take full width when set to \texttt{display: block}.
\end{solution}

\endgradingrange{csawebstatic}

